\documentclass[11pt]{article}
\usepackage[a4paper,margin=27mm]{geometry}
\usepackage{amsmath,amssymb,mathtools}
\usepackage{fontspec,unicode-math}
\setmainfont{Latin Modern Roman}
\setmathfont{Latin Modern Math}
\usepackage{xurl,hyperref}
\hypersetup{colorlinks=true,urlcolor=blue,linkcolor=blue}
\newcommand{\sref}[2]{\hyperlink{source:#1:#2}{[S#2]}}
\newcommand{\cataloguescope}[1]{\paragraph{Mathematical question.} #1}
\begin{document}
\section{Reverse Mathematics of Fraïssé's Conjecture}\label{statement}

\subsection{Definitions and mathematical statement}

\paragraph{Shared notation.}$\mathbb N=\{1,2,\ldots\}$, and $\mathbb Z,\mathbb Q,\mathbb R,\mathbb C$ denote the integers, rationals, reals and complex numbers. Any other conventions are specified locally.


Work over $\mathsf{RCA}_0$, the subsystem of second-order arithmetic with $\Delta^0_1$ comprehension and $\Sigma^0_1$ induction; all structures and orders are countable and coded by sets of integers. $\mathsf{ACA}_0$ adds comprehension for every arithmetical formula (set parameters allowed), and $\mathsf{ATR}_0$ adds arithmetical transfinite recursion along every coded well-order. Reverse-mathematical equivalence means provability of both implications over this base. For linear orders $L,K$, write $L\preccurlyeq K$ if some injective map $e:L\to K$ preserves the order. Let $\mathsf{FRA}$ say that every sequence $(L_i)_{i\in\mathbb N}$ of countable linear orders has $i<j$ with $L_i\preccurlyeq L_j$. Let $\Pi^1_1$-$\mathsf{CA}_0$ denote comprehension for $\Pi^1_1$ formulas, including set parameters.
\paragraph{Statement recorded in the supplied source.}
The historical question has two parts: is $\mathsf{FRA}$ equivalent to $\mathsf{ATR}_0$ over $\mathsf{RCA}_0$, and is $\mathsf{FRA}$ provable in $\Pi^1_1$-$\mathsf{CA}_0$? The second part has an affirmative answer in the later Theorem 1.1 cited below. \sref{statement}{3} \sref{statement}{2}
\subsection{Short English statement}
Determine whether Fraïssé's embeddability principle has exactly the strength of arithmetical transfinite recursion; its historical question about a stronger comprehension upper bound was answered affirmatively.
\subsection{Sources}
\begin{description}
\item[\hypertarget{source:statement:1}{S1}] ULAM.AI and the UnsolvedMath Contributors, UnsolvedMath: A Curated Collection of Open Mathematics Problems, catalogue entry. CC BY 4.0. \url{https://huggingface.co/datasets/ulamai/UnsolvedMath}.
\item[\hypertarget{source:statement:2}{S2}] Antonio Montalbán, Open Questions in Reverse Mathematics, author preprint (17 August 2010), Question 21, Section 4.2, p. 9. \url{https://math.berkeley.edu/~antonio/papers/questionsRM.pdf}.
\item[\hypertarget{source:statement:3}{S3}] Antonio Montalbán, Fraïssé's conjecture in $\Pi^1_1$-comprehension, author preprint (compiled 2017), Theorem 1.1. \url{https://math.berkeley.edu/~antonio/papers/FraisseDeltaBQO.pdf}.
\end{description}
\label{end:statement}
\end{document}
